\section{Category Theory}

\begin{definition}[Category Universe]
A category $\mathcal{C} \in \Cat$ consists of a collection of objects $\Ob(\mathcal{C})$, a collection of morphisms $\Mor(\mathcal{C})$, an identity assignment $\id_A \in \Hom(A, A)$ for each $A \in \Ob(\mathcal{C})$, and associative composition.
\end{definition}

\begin{definition}[Functors and Natural Transformations]
Let $F, G \colon \mathcal{C} \to \mathcal{D}$ be functors with opposite category $\mathcal{C}^{\op}$. The morphism actions are denoted $\FMor$ and $\GMor$. A natural transformation $\eta \in \Nat(F, G)$ assigns to each object $A \in \Ob(\mathcal{C})$ a component $\eta_A \in \Hom(F(A), G(A))$.
\end{definition}

% Commutative square or diagram figure fixture
\begin{figure}[H]
\centering
\setlength{\unitlength}{1.2cm}
\begin{picture}(4,3)
  \put(0.8,2.5){\makebox(0,0){$A$}}
  \put(3.2,2.5){\makebox(0,0){$B$}}
  \put(0.8,0.5){\makebox(0,0){$C$}}
  \put(3.2,0.5){\makebox(0,0){$D$}}

  \put(1.2,2.5){\vector(1,0){1.6}}
  \put(2.0,2.8){\makebox(0,0)[b]{$f$}}

  \put(1.2,0.5){\vector(1,0){1.6}}
  \put(2.0,0.2){\makebox(0,0)[t]{$g$}}

  \put(0.8,2.1){\vector(0,-1){1.2}}
  \put(0.5,1.5){\makebox(0,0)[r]{$h$}}

  \put(3.2,2.1){\vector(0,-1){1.2}}
  \put(3.5,1.5){\makebox(0,0)[l]{$k$}}
\end{picture}
\caption{Morphism composition diagram demonstrating $k \circ f = g \circ h$.}
\label{fig:commutative_square}
\end{figure}


\begin{lemma}[Hom Bifunctor]
The hom-set assignment induces a lifted morphism functor $\HomMor \colon \mathcal{C}^{\op} \times \mathcal{C} \to \Set$.
\end{lemma}

\myprf
Verification of identity and composition preservation follows from the diagram in Figure~\ref{fig:commutative_square}.
\myqed
